\section{Introduction}

Prescribing the probability that each edge belongs to a random perfect
matching does not prescribe the distribution of the whole matching.
The largest possible entropy of such a distribution measures how many
matching choices can coexist with those edge probabilities. We bound
this entropy by explicit sums of one-coordinate functions, at every
feasible mean. The result yields lower counts from uniform marginals
and additive estimates for weighted partition functions.

Throughout, a graph is finite, undirected and loopless. Parallel edges
are distinguished both as coordinates and as choices in a matching.
A perfect matching covers each vertex by exactly one edge. Write
$\mathcal M(G)$ for the set of perfect matchings,
$N(G)=|\mathcal M(G)|$, and
\[
 P(G)=\conv\{\mathbf1_M:M\in\mathcal M(G)\}\subseteq\R^{E(G)}
\]
for their polytope. Thus a point $x\in P(G)$ is exactly a feasible
vector of edge probabilities. For a probability law $p$ on the matching
family, put $\Ent(p)=-\sum_M p_M\log p_M$, and define
\begin{align}
 H(x)&=\max\left\{\Ent(p):p\text{ is a probability law on }\mathcal M(G),\
                   \sum_Mp_M\mathbf1_M=x\right\},\label{intro:H}\\
 F(x)&=-\sum_{e\in E(G)}x_e\log x_e,\qquad
 B(x)=-\sum_{e\in E(G)}(1-x_e)\log(1-x_e).
 \label{intro:FB}
\end{align}
All unqualified logarithms are natural, and $0\log0=0$.
The maximum in Equation~\eqref{intro:H} is attained by compactness
of the feasible set of laws. No connectedness or bipartiteness is assumed.

\begin{theorem}[Pointwise entropy bounds]\label{ent:main}
Let $G$ be a loopless multigraph on $2m\ge2$ vertices with at least one
perfect matching. For every $x\in P(G)$,
\begin{equation}\label{intro:main-bound}
                  F(x)-\left(2-\frac2m\right)B(x)
                  \le H(x)\le F(x).
\end{equation}
\end{theorem}

In particular, $H\ge F-2B$ throughout the polytope, including its
boundary. When $m=1$, each matching is one of the distinguished edges
between two vertices, and the bounds give the exact identity $H=F$.
For the empty vertex set the unique empty matching gives $F=B=H=0$;
the formula involving $1/m$ is not used.

The upper bound $H\le F$ is a direct instance of Shearer's entropy
inequality~\citep{ChungGrahamFranklShearer1986}: each edge occurs in
two vertex stars, and a matching chooses exactly one edge at each
vertex. Our contribution to Equation~\eqref{intro:main-bound} is its
lower comparison. The coefficient there is obtained from a geometric
rank estimate followed by a covariance calculation, with no assumption
that the mean is interior to the full polytope.

Edmonds' matching-polytope theorem~\citep{Edmonds1965} describes feasible
matching means by degree equations, nonnegativity and odd-cut
inequalities. The odd cuts distinguish general graphs from bipartite
graphs and are also central to the entropy bound here. Bipartite graphs
admit an order-independent coefficient-one comparison:
Schrijver's inequality~\citep{Schrijver1998Counting} and the weighted Bethe
inequality of Gurvits~\citep[expanded version, Theorem~2.2]{Gurvits2011SchrijverBethe}, together with
entropy duality, imply $H(x)\ge F(x)-B(x)$. Their permanent setting is
essential to this comparison. In Section~\ref{ex:sec-triangle-expansions} we give a
simple eight-vertex nonbipartite graph with a feasible mean for which
$H<F-B$.

\subsection{Entropy and structural consequences}

Anari, Oveis Gharan and Vinzant~\citep[FOCS version, Conjecture~6]{AnariOveisGharanVinzant2018}
proposed a convex-optimization framework for approximate counting
and asked whether the maximum binary-coordinate entropy of a perfect
matching mean exceeds $\log N(G)$ by only a linear function of the
number of vertices. In our notation their objective is $F+B$.
Theorem~\ref{ent:main} gives the explicit positive answer
\[
 \log N(G)\le\max_{x\in P(G)}(F(x)+B(x))
                       \le\log N(G)+3m-2.
\]
The pointwise statement also controls arbitrary finite edge weights,
through the Gibbs variational formula. The proposed linear-error
estimate is therefore an optimized consequence of the prescribed-mean
inequality, rather than its definition. Corollary~\ref{cons:entropy-count} proves the optimized bounds, and
Section~\ref{cons:weighted-section} derives the weighted form.

If $G$ contains $k\ge2$ edge-disjoint perfect matchings, their uniform
mixture has marginal $1/k$ on their union. Substituting that mean into
Theorem~\ref{ent:main} gives
\[
 N(G)\ge k^m(1-1/k)^{2(m-1)(k-1)}
       \ge e^{-2(m-1)}k^m.
\]
The same estimate holds for a $k$-regular graph whose every odd cut has
at least $k$ edges, because its uniform vector belongs to $P(G)$.
These consequences and their endpoint cases are proved in
Section~\ref{cons:section}.

For structural counts, the scale of the degree matters. Esperet,
Kardo\v{s}, King, Kr\'a\v{l} and Norine~\citep{EsperetEtAl2011}
proved an exponential lower bound for every bridgeless cubic graph.
Abdi, Cornu\'ejols, Dadush and Dalirrooyfard~\citep[Corollary~6.6]{AbdiCornuejolsDadushDalirrooyfard2026}
proved exponential lower bounds for all regular graphs of degree at
least four satisfying the corresponding odd-cut condition, including
matchings crossing every minimum odd cut exactly once.
Ebrahimnejad, Nagda and Oveis Gharan~\citep{EbrahimnejadEtAl2021}
gave quantitative lower counts for regular nonbipartite graphs under
an additional spectral-expansion hypothesis. The estimate
$(k/e^2)^m$ obtained here has a base growing with the degree $k$.
At small fixed degree these earlier exponential results can be
stronger: our displayed homogeneous base is greater than one only
when $k\ge8$. On the bipartite subclass Schrijver's bound has the
stronger asymptotic base $k/e$ per matched pair. Our structural
statements concern the explicit odd-cut condition; merely requiring
every edge to belong to a perfect matching is a different hypothesis.
Esperet et al. formulate the growing-degree question on the broader
matching-covered class \citep[corrected full version, Conjecture~24]{EsperetEtAl2011}.
Our regular odd-cut corollary also establishes the even-order connectivity
formulation of \citet[arXiv version, Conjecture~1.1]{EbrahimnejadEtAl2021}.\footnote{The
cited arXiv version writes the degree as $d$ but prints an undeclared
$k-1$ in the connectivity condition. We read it as $d-1$, consistently
with its Section~1.2.}

\subsection{Face dimension and the entropy argument}

For a graph on $2m\ge2$ vertices and $x\in P(G)$, let $F_x$ be its
minimal face, the smallest face of
$P(G)$ containing $x$, and let $\supp x=\{e:x_e>0\}$.
The geometric input is the sharp estimate
\[
                  |\supp x|-\dim F_x\le3m-2.
\]
Edmonds' theorem identifies the left side with the rank of the tight
odd-cut equations on the supported coordinates; singleton cuts supply
the degree equations. Abdi et al.~\citep[Lemmas~6.4--6.5]{AbdiCornuejolsDadushDalirrooyfard2026}
bound this rank by $3m-1$ using odd laminar families.
Their cubic staircase examples already have minimum-odd-cut rank
$|E|-2$, attaining $3m-2$
\citep[discussion after Corollary~6.6 and Figure~2]{AbdiCornuejolsDadushDalirrooyfard2026}.

The dimension of the whole perfect-matching polytope was characterized
by Naddef~\citep{Naddef1982} and by Edmonds, Pulleyblank and
Lov\'asz~\citep{EdmondsPulleyblankLovasz1982}, whose brick decomposition
uses cuts crossed once by every perfect matching. Here a cut need only
be tight at the prescribed mean: it is crossed once by every matching
in that minimal face, but possibly not by every matching of the graph.
Section~\ref{geom:sec-polytope} proves the supported bound by contracting
such a cut. A direction in the parent face can be assembled from
directions in the two contracted faces when their crossing-edge
coordinates agree. Each crossing-coordinate sum is zero, by the degree
equation at the contracted vertex. Agreement therefore imposes at most
$|C|-1$ independent conditions for a cut $C$, rather than $|C|$.
This saved condition closes an induction with bound $3m-2$.
Appendix~\ref{a:section} also gives the laminar proof. It uses the standard
face-restricted uncrossing argument
\citep[full version, Lemma~2.2 and Appendix~A]{SvenssonTarnawski2017}
and retains the dependence between complementary root children in the
laminar tree.

The analytic step uses the classical exponential form of a
maximum-entropy law under mean constraints
\citep{Jaynes1957,SinghVishnoi2014}. Its log-partition Hessian is the
covariance matrix, as in standard exponential-family theory
\citep{WainwrightJordan2008}. We restrict the law to the atoms of $F_x$.
Its covariance then has image equal to the direction space of that face,
even when it is singular in the ambient edge space.

Suppose that $F-cB-H$ has a positive maximum. Removing coordinates fixed
at zero or one leaves strictly fractional means and preserves this
entropy deficit. Exponential tilts within the resulting minimal face
supply smooth curves through the maximum. First-order stationarity
cancels the acceleration term in their second derivatives. Summing the
second-derivative inequalities over positive covariance eigenvalues
bounds the face dimension above, contradicting the geometric lower bound
whenever $c>2-2/m$ for the original graph. Passing to the endpoint proves
Theorem~\ref{ent:main}. Section~\ref{sec:entropy} carries this out within
the face, without inverting an ambient covariance kernel or differentiating
across faces.

Section~\ref{ex:sec-triangle-expansions} explains both geometric sharpness
and the obstruction to coefficient one. We use the triangle expansions
underlying the Klee graphs of
Cygan, Pilipczuk and \v{S}krekovski~\citep{CyganPilipczukSkrekovski2013}
to construct a two-dimensional minimal face on every positive even
order. At uniform edge marginals, a matching law is forced onto its
three vertices even when the graph has further matchings. These examples
attain the geometric bound and give the coefficient-one counterexample;
they do not establish optimality of the coefficient two.

\subsection{Deterministic counting with binary multiplicities}

For the algorithmic problem, the input specifies a vertex set and
nonnegative integer multiplicities $k_{ij}$ in binary for unordered pairs
of distinct vertices; omitted pairs have multiplicity zero. Its count is
\[
 N(G)=\sum_M\prod_{\{i,j\}\in M} k_{ij},
\]
where $M$ runs over pairings of the vertex set. The input can also list
parallel edges explicitly. The algorithm uses the binary multiplicities
without expanding them into labelled copies.

\begin{theorem}[Deterministic approximate counting]
\label{thm:main}\label{thm:algorithm-guarantee}
There is one uniform deterministic algorithm which, on a graph on $n$
vertices with nonnegative binary integer pair multiplicities, returns an
integer $A\ge0$ satisfying
\[
                         N(G)/512^n\le A\le N(G).
\]
Its running time and output length are polynomial in the input bit length.
In particular, $A=0$ exactly when $N(G)=0$. The guarantee includes a compact
input consisting of binary $n$ and a list of positive pair records.
\end{theorem}

The factor $512^n$ permits a complete implementation with integer and
rational arithmetic and a fixed polynomial number of steps. The degree
of this polynomial is large; the theorem gives a complexity guarantee,
not a practical running-time benchmark. Appendix~\ref{c:app-loops} treats
the different convention that a loop covers its one incident vertex. It
returns a rational approximation within factor $2^{18n}$, including when
the original vertex count is odd.

\paragraph{Approximation history and computational models.}

Exact counting is $\#\mathrm P$-complete even for bipartite perfect
matchings, through the permanent of a zero--one matrix
\citep[Theorem~1]{Valiant1979}. For an $r\times r$ nonnegative matrix,
Linial, Samorodnitsky and Wigderson give an $e^r$-factor deterministic
strongly polynomial arithmetic algorithm
\citep[author version, Theorem~1.1]{LinialSamorodnitskyWigderson2000}.
Jerrum, Sinclair and Vigoda give a fully polynomial randomized
approximation scheme for the nonnegative permanent
\citep{JerrumSinclairVigoda2004}. For nonnegative rational matrices, Kudria, Luo and Majid give certified
rational endpoints in polynomial bit time with logarithmic width
$O(r(\log\log r)^2/\log r)$
\citep[Theorem~1.1]{KudriaLuoMajid2026}, while Dong and Jain give factor
$(1+\varepsilon)^r$ for every fixed $0<\varepsilon\le1$ in a strongly
polynomial arithmetic model
\citep[Theorem~1.1 and Section~1.1]{DongJain2026}.
These permanent results concern bipartite graphs on $2r$ vertices.

The analogous function for a general graph is the \emph{hafnian}: the
sum of products over pairings of the indices of a symmetric matrix.
Barvinok describes a randomized exponential-factor approximation for
nonnegative hafnians \citep[author preprint, Section~7]{Barvinok1999}.
For fixed $\delta>0$ and any $0<\varepsilon<1$, his later theorem approximates
$\ln\operatorname{haf}(A)$ to additive error $\varepsilon$ using
$r^{O_\delta(\log r+\log(1/\varepsilon))}$ arithmetic operations when the entries of the
symmetric $2r\times2r$ matrix lie in $[\delta,1]$
\citep[Theorem~2.1]{Barvinok2017}. The diagonal may be filled arbitrarily
in this interval, since a hafnian uses no diagonal entry.
Yi gives a deterministic fully polynomial approximation scheme for
rational hafnian inputs of even order, with zero diagonal, when nonzero
entries lie in $[\theta,1]$ and the support has minimum degree at least
$(1/2+\gamma)$ times its order, for fixed $0<\theta\le1$ and
$0<\gamma<1/2$ \citep[Theorem~1.1(a)]{Yi2026}. These hypotheses differ from arbitrary supports and binary
multiplicities in Theorem~\ref{thm:main}. A separate manuscript treats
fully polynomial randomized relative approximation for arbitrary simple
graphs \citep[Theorem~1.1]{OpenAIMatchingFPRAS2026}.

The construction follows the entropy-optimization approach of
Anari, Oveis Gharan and Vinzant~\citep[FOCS version, Section~I.D]{AnariOveisGharanVinzant2018},
with a rational optimizer and all bit bounds given explicitly here.
Section~\ref{cons:weighted-section} establishes the binary-log weighted
interface. Sections~\ref{d:sec-objective} and~\ref{d:sec-optimization}
first complete the input to $K_n$, assigning tiny positive dyadic
weights to absent pairs. This keeps the optimization polytope nonempty
and separates zero from positive original counts. They then replace
the weighted entropy objective by an integer-slope dyadic interpolant
and an exact penalty on a box. On $K_n$, correcting degrees and then
mixing with the uniform feasible point repairs constraint violations. The resulting fixed-grid
iteration needs only an active violated row. Section~\ref{d:sec-odd-cuts}
finds that row by a parity-preserving cut recursion, and
Section~\ref{d:sec-algorithm} proves the count, zero-detection and
bit-complexity guarantees. Minimum odd cuts as matching-constraint
separators go back to Padberg and Rao~\citep{PadbergRao1982}.
The exponential matching laws used in the entropy proof remain analytic
objects; this algorithm does not query them.

The appendices retain alternative relations between face geometry and
entropy. Appendix~\ref{a:section} gives the laminar and quotient-map
arguments, then the direct coefficient-three and coefficient-four
covariance comparisons and their counting consequences.
Appendix~\ref{b:section} proves the nonlinear entropy-defect bound by
normalizing positive marginals and taking the inverse covariance on
the face's direction space; Appendix~\ref{a:section} also gives its
short scalar deduction. Appendix~\ref{c:app-loops} proves the
singleton-loop reduction and its rational bit complexity.
